\section{HC and Lepton Mass Readout: An Identification Test}
\label{v480:hc-mass-identification}

The sufficient construction in Proposition~\ref{v480:koide-hc-mixing}
is not yet an identification with the Companion's HC. This section tests
that identification against the actual inherited definitions. Its
negative results concern specified readout routes, not every possible
nonlinear, coupled or symmetry-breaking realization of the framework.

\subsection{What the Existing Definitions Supply}
\label{v480:hc-mass-source-map}

The inherited decomposition is $C=(I+J)/2$, $P_-=(I-J)/2$, with
$CP_-=0$; see Sections~\ref{sec:central} and~\ref{sec:stringlqg}.
The fermionic input lane is declared $J$-odd, while $C$ selects the
compensated photon/fix output. The refined mass discussion instead
assigns massive readout to localization remaining relative to that
output (Section~\ref{subsec:mass-null-axis}) and names Higgs mediation
as its interface (Definition~\ref{def:higgs-mediation}).

Proposition~\ref{prop:spectral-to-mass-realization-map} names
$\mathcal R_m:\mathcal S_{\rm spec}\to\mathcal S_{\rm mass}$,
but its argument does not provide a numerical mass functional,
generation-resolved coupling, or action on three principal mass roots.
The inherited remark ``Mass spectrum beyond Koide'' explicitly leaves
the Yukawa mechanism undeveloped. The symbolic effective operator in
Theorem~\ref{thm:minimal-mmhc-field-equation} does not supply these
missing matrix elements or a lepton spectral assignment either.

Two different triples must not be identified by their cardinality:
$(K_1,K_2,K_3)$ denotes matter/tachyon/photon sector readouts,
whereas $(e,\mu,\tau)$ denotes three generations within the lepton
readout. The generation discussion in Theorem~\ref{thm:three-generations} does not
construct a map equating their components or their scalar products.

There is also a scalar-field issue. The inherited unitary-equivalence
argument calls $J$ unitary, whereas the proof of
Theorem~\ref{thm:gauge-invariance} calls it anti-unitary. An
anti-unitary involution and $(I+J)/2$ are generally only real-linear.
The following operator tests therefore explicitly use a real Hilbert
realization, where $J$ is an orthogonal involution. They do not silently
assert a complex-linear three-dimensional representation of either
inherited convention.

\subsection{The Exact Linear HC-to-Root Interface}
\label{v480:hc-mass-interface}

\begin{proposition}[A sufficient norm-controlled descent]
\label{v480:hc-mass-descent}
Let $H$ be a real Hilbert space, $J$ an orthogonal involution, and
$B:H\to\mathbb R^3$ a bounded coisometry, $BB^*=I_3$. If
$\ker B$ is $J$-invariant, then
\[
 J_m=BJB^*,\qquad J_m^*=J_m,\qquad J_m^2=I_3,
 \qquad BC=C_mB,\quad C_m=(I_3+J_m)/2.
\]
Conversely, any relation $BJ=J_mB$ requires $\ker B$ to be
$J$-invariant. The descended operator is unique because $B$ is onto.
\end{proposition}
\begin{proof}
Put $E=B^*B$, the orthogonal projection onto $(\ker B)^\perp$.
Since $J$ is self-adjoint and preserves $\ker B$, it also preserves
its orthogonal complement, so $JE=EJ$. Thus $J_m^*=J_m$ and
$J_m^2=BJEJB^*=BEB^*=I_3$. Also
$J_mB=BJE=BJ$, since $J$ preserves $\ker B$ and $B$ kills it.
The identity for $C$ follows by averaging. Conversely,
$Bk=0$ implies $BJk=J_mBk=0$. Surjectivity gives uniqueness.
\end{proof}

Coisometry is a sufficient metric-control assumption, not a consequence
of merely naming a readout. Without invariant kernel, compression
$BJB^*$ need not even be an involution. A general invertible change of
coordinates gives an involution $BJB^{-1}$, but need not preserve the
Euclidean scalar product used in Koide. A different root-space metric
cannot be substituted without re-deriving the physical ratio.

If an independently supplied seed $f_0$ satisfies $Bf_0=e$ and the
\emph{physical} root readout is proved to be
\[
 x=\alpha BCf_0=\alpha C_m e,\qquad
 e=(1,1,1)/\sqrt3,\quad \alpha>0,
\]
with nonzero nonnegative components, the earlier proposition applies.
Here $\alpha$ carries square-root mass units for a dimensionless seed.
Its missing condition can now be stated on the original carrier:
\[
 \langle e,J_me\rangle
   =\langle B^*e,JB^*e\rangle=0.
\]
The paper presently supplies neither this $B$, this seed/root law,
nor that independently justified orthogonality. The displayed
interface is one linear route to investigate, not a requirement that
every physical mass readout be linear. In particular, a globally
nonnegative linear root map on all of $H$ would have to vanish, since
both $Bf$ and $B(-f)=-Bf$ would be nonnegative. A nonzero physical
linear root law must be restricted to an admitted state cone.

\subsection{Obstructions to the Direct Identification}
\label{v480:hc-mass-obstructions}

\begin{proposition}[The pure HC output loses the odd input]
\label{v480:hc-mass-odd-obstruction}
If $Jf=-f$, then $Cf=0$. Consequently any homogeneous linear
root readout $x=BCf$ is zero. Any readout $F(Cf)$ with $F(0)=0$
has the same obstruction on this input lane. If instead a seed is
$J$-fixed and the descended linear route has $Bf_0=e$, then
$J_me=e$ and its nonzero root output has $K=1/3$.
\end{proposition}
\begin{proof}
The first statement follows from $Cf=(f+Jf)/2=0$.
For the second, $J_me=BJf_0=Bf_0=e$, so $C_me=e$ and
all three roots are equal.
\end{proof}

This rules out direct identification of the declared odd fermionic
input with a positive root vector obtained only from its compensated
output. It does not rule out a mass readout of the residual input or
of a coupled state. Nor does quaternionic orthogonality establish the
needed root correlation: for an orthogonal projection $C$ one has
$\langle Cf,f\rangle=\|Cf\|^2$. Its nonzero output cannot be
orthogonal to its own input in that same Hilbert metric. The inherited
quaternion-sector wording in Proposition~\ref{prop:compensator} must
not be used as the proof of $\langle e,J_me\rangle=0$.

\begin{proposition}[Equivariant readout cannot split a fixed input]
\label{v480:hc-mass-equivariant-obstruction}
Let $U$ preserve an admitted state set and let $P$ cyclically permute
the three root coordinates. For any readout $F$ satisfying
$F(Uf)=PF(f)$, an input $Uf_0=f_0$ has
$F(f_0)=a(1,1,1)$. If this is a nonzero nonnegative root vector,
then $K=1/3$. No linearity assumption on $F$ is needed.
\end{proposition}
\begin{proof}
Equivariance gives $PF(f_0)=F(Uf_0)=F(f_0)$. The fixed vectors
of the three-cycle are exactly the constant triples.
\end{proof}

A common frame parameter does not evade this result if the input
remains fixed and the readout remains equivariant at each parameter
value. An asymmetric input, a selected symmetry-breaking state or
an explicitly non-equivariant coupling is needed. The generation
action $U$ and the matter/tachyon involution $J$ are distinct data.

\begin{proposition}[Closed-triole amplitudes do not yield Koide by norms]
\label{v480:hc-mass-triole-bound}
Let $z_0+z_1+z_2=0$ in a normed vector space, not all $z_i=0$.
Under the specific identification $m_i=\kappa\|z_i\|^2$ with a
common $\kappa>0$, the principal-root ratio satisfies
\[
 \frac13\leq K\leq\frac12.
\]
In particular it cannot equal $2/3$.
\end{proposition}
\begin{proof}
Put $r_i=\|z_i\|$ and $s=\sum_i r_i>0$. The triangle inequality
gives $r_i\leq s-r_i$, hence $r_i\leq s/2$. Therefore
$\sum_i r_i^2\leq(s/2)\sum_i r_i=s^2/2$.
Cauchy--Schwarz gives the lower bound, and the common mass scale
cancels from $K$.
\end{proof}

This excludes a natural but additional identification, not every
generation readout of a triole. Unequal generation couplings or a
different nonlinear realization need not obey this bound; they are
precisely additional data that must be constructed rather than read
off from $z_0+z_1+z_2=0$ alone.

\begin{proposition}[An unrestricted fitted HC shape predicts no roots]
\label{v480:hc-mass-fit-freedom}
For every nonzero nonnegative triple $x$, there exist an orthogonal
involution $J_x$ and $\alpha_x>0$ such that
$x=\alpha_x(I+J_x)e/2$. Thus existence of an HC-shaped operator
chosen after the masses is not a prediction of Koide or a hierarchy.
\end{proposition}
\begin{proof}
Set $C_x=xx^{\mathsf T}/\|x\|^2$ and $J_x=2C_x-I_3$.
This is an orthogonal involution. Since $s=\sum_i x_i>0$,
$C_xe=sx/(\sqrt3\|x\|^2)$, so choose
$\alpha_x=\sqrt3\|x\|^2/s$.
\end{proof}

\subsection{A Residual Coupling Test Compatible With the Mass Language}
\label{v480:hc-mass-residual-test}

The refined massive-branch definition suggests testing the component
that $C$ removes, rather than identifying the masses with $Cf$.
The following is an explicit \emph{additional effective ansatz}, not
a recovered Yukawa interaction. Put $H_-=\ker(J+I)$ and normalize
the declared states dimensionlessly. Fix a mass scale $\kappa>0$ and a dimensionless
bounded linear coupling $D_\beta:H_-\to V$ to a declared real
Hilbert output space. For three admitted generation states define
\[
 u_i=P_-\psi_i,\qquad
 m_i(\beta)=\kappa\|D_\beta u_i\|^2,\qquad
 x_i=\sqrt\kappa\|D_\beta u_i\|.
\]
The quadratic form is positive, is zero on a fully fixed input,
and is unchanged when $\psi_i$ is replaced by $J\psi_i$, since
$P_-J\psi_i=-P_-\psi_i$. Unlike a pure $BC\psi_i$ law, it can
be nonzero on the declared fermionic input lane. It assigns the
same mass to $J$-dual inputs; a different dual-sector mass law
would require a different frame-labelled coupling or readout.

For $q_i=\|D_\beta u_i\|^2$, not all zero, Koide still requires
the independent equation
\[
 \sum_i q_i=\frac23\left(\sum_i\sqrt{q_i}\right)^2.
\]
Neither positivity nor the residual decomposition proves it.
In particular, if all $\|u_i\|=1$ and
$D_\beta^*D_\beta=\beta^2 I$ on their span, all masses are
$\kappa\beta^2$ and $K=1/3$ for $\beta\ne0$. Thus the isotropic
frame term $\beta^2\langle f,f\rangle$ in
Definition~\ref{def:lagrangian}, used alone on equally normalized
generations, does not generate a hierarchy. A common positive
rescaling leaves $K$ unchanged. At $\beta=0$ in this particular
ansatz all masses vanish and $K$ is undefined, not a nonzero
degenerate Koide limit.

\paragraph{Identification status and the next concrete object.}
No native HC-to-lepton-root identification is established by the
current definitions. The pure odd-input HC route and the common-norm
closed-triole route are excluded under their stated assumptions.
The remaining constructive target suggested by the existing
mass/Higgs language is an independently specified generation-resolved
coupling $D_\beta$ and generation states $\psi_i$, or an explicit
alternative mass realization. Their domains, generation action,
positivity, physical units and relation to the effective field
operator must be supplied before testing the equal-norm/Koide law,
the phase and the absolute scale. Assigning entries from the observed
masses would only fit the missing map. The inherited assertions that
historical OP~3 is closed, and P2 is confirmed, do not discharge this
obligation; their physical derivation remains unestablished here.
